\section{The complete irreducibility classification in A327315}
\label{sec:a327315}

\subsection{The literal source statement}

Let $P_n$ be the numerator of the $n$th derivative of
\[
 f(X)=\frac{X-2}{X^2-X+1},
\]
divided by $n!$, with the rational function reduced and its denominator
monic.  OEIS A327315~\cite{A327315} records these rows from $n=0$ and conjectures that
their irreducible indices are exactly A069353~\cite{A069353}, the numbers $2^a3^b-1$.
The entry attributes the conjecture to Clark Kimberling in November 2019.
Its title says degree $n$, but the executable definition and displayed
rows have degree $n+1$, beginning with $P_0=X-2$.  We follow those
unambiguous row and derivative indices.

\begin{theorem}[All derivative orders]
\label{thm:a327315}
For every integer $n\ge0$,
\[
 P_n\text{ is irreducible over }\mathbb Q
 \quad\Longleftrightarrow\quad
 n+1=2^a3^b\text{ for some }a,b\ge0.
\]
The same criterion holds for the unnormalized reduced numerator $n!P_n$.
\end{theorem}

There is one classification here.  The stronger irreducibility over a
quadratic field, the explicit factors in the excluded degrees, and the
source's degree correction are not counted as additional conjectures.

\subsection{Recurrence and exact reduced numerator}

Put $A=1-2X$ and $D=1-X+X^2$, and define
\[
 S_0=-1,\qquad S_1=X-2,\qquad
 S_{N+2}=A S_{N+1}-D S_N.
\]
The top-coefficient recurrence has initial values $-1,1$ and satisfies
$c_{N+2}=-2c_{N+1}-c_N$.  Consequently
\[
 \deg S_N=N,\qquad \operatorname{lc}(S_N)=(-1)^{N+1}.
\]
Differentiating the recurrence and using two-step induction gives
\[
 S_N'=-N S_{N-1}\quad(N\ge1).
\]
Combining this with the recurrence yields, for all $N\ge0$,
\begin{equation}
\label{eq:a327315-differential}
 D S_N'+N A S_N=N S_{N+1}.
\end{equation}
Since $D'=-A$, the quotient rule and
\eqref{eq:a327315-differential} prove inductively that
\begin{equation}
\label{eq:a327315-derivative}
 f^{(n)}(X)=\frac{n!S_{n+1}(X)}{D(X)^{n+1}}.
\end{equation}

The identities
\[
 4D-A^2=3,\qquad D-(X+1)(X-2)=3
\]
give Bezout certificates over $\mathbb Q$ for $(A,D)$ and $(S_1,D)$.
The polynomial $S_0=-1$ is a unit.  Modulo $D$, the recurrence says
$S_{N+2}\equiv A S_{N+1}$, so induction proves
$\gcd(S_N,D)=1$ for every $N$.  Thus the fraction in
\eqref{eq:a327315-derivative} has no polynomial cancellation.  Its
denominator is monic, and $n!$ is a nonzero rational unit.  Therefore its
canonical numerator is exactly $n!S_{n+1}$, and
\begin{equation}
\label{eq:a327315-source}
 P_n=S_{n+1}.
\end{equation}
This identifies the recurrence with the original reduced-numerator
definition, including $n=0$.

\subsection{The quadratic field and fractional-linear form}

Let $K=\mathbb Q(u)$ with $u^2-u+1=0$, and set $v=1-u$.
The polynomial $Y^2-Y+1=\Phi_6(Y)$ is irreducible over $\mathbb Q$, so
$[K:\mathbb Q]=2$.  Direct algebra gives
\[
 u+v=uv=1,\quad u^2=-v,\quad v^2=-u,\quad
 u^3=v^3=-1,\quad u^6=v^6=1.
\]
Also $u\ne v$, and neither $u$ nor $v$ is zero or one.
With $\alpha=u-X$ and $\beta=v-X$, we have
$\alpha+\beta=A$ and $\alpha\beta=D$.  The initial values and
recurrence therefore prove
\begin{equation}
\label{eq:a327315-closed}
 S_N=-v\alpha^N-u\beta^N
 =(-1)^{N+1}v\big((X-u)^N-v(X-v)^N\big).
\end{equation}

We use the following standard fractional-linear observation.  Over any
field, if $a\ne b$, $c\ne1$, and $Y^N-c$ is irreducible, then
\[
 (X-a)^N-c(X-b)^N
\]
is irreducible of degree $N$.  Indeed, if $x$ is a root of an irreducible
factor, then $x\ne b$ and
$y=(x-a)/(x-b)$ satisfies $y^N=c$.  Thus the extension containing $x$
has degree at least $N$, while the displayed polynomial has degree $N$.
Every irreducible factor must have full degree.  The Lean helper proves
this argument directly with an adjoint-root field.

\subsection{Irreducibility in every allowed degree}

\begin{lemma}
\label{lem:a327315-cyclotomic}
For $N=2^a3^b$, with $a,b\ge0$,
\[
 Y^{2N}-Y^N+1=\Phi_{6N}(Y).
\]
\end{lemma}
\begin{proof}
The initial identity is $\Phi_6(Y)=Y^2-Y+1$.  If a prime $p$ divides
$r$, the classical cyclotomic expansion formula is
$\Phi_r(Y^p)=\Phi_{rp}(Y)$.  Starting at $r=6$, apply this formula
$a$ times with $p=2$ and $b$ times with $p=3$.  Both required divisibilities
hold at every step, and the left side becomes the stated trinomial.
\end{proof}

\begin{lemma}
\label{lem:a327315-binomial}
If $N=2^a3^b$, then $Y^N-v$ is irreducible over $K$.
\end{lemma}
\begin{proof}
Choose an irreducible factor $q\in K[Y]$, of degree $d$, and let $y$ be
its root in $E=K[Y]/(q)$.  Then $y^N=v$, and hence
\[
 y^{2N}-y^N+1=v^2-v+1=0.
\]
By Lemma~\ref{lem:a327315-cyclotomic}, this is the irreducible polynomial
$\Phi_{6N}$ over $\mathbb Q$, of degree $2N$.  It is therefore the
minimal polynomial of $y$ over $\mathbb Q$.  The tower law gives
\[
 2N\le[E:\mathbb Q]=[K:\mathbb Q][E:K]=2d.
\]
Since $q$ divides a degree-$N$ polynomial, $d\le N$.  Thus $d=N$, proving
the assertion.
\end{proof}

Equation~\eqref{eq:a327315-closed} and the fractional-linear observation
now prove that $S_N$ is irreducible over $K$ for every $N=2^a3^b$.
It follows that $S_N$ is irreducible over $\mathbb Q$ as well: a proper
rational factorization would remain a proper factorization over $K$.

\subsection{Explicit proper factors in every excluded degree}

Let $N>0$ have a prime divisor $p\notin\{2,3\}$, and write $N=pm$.
Then $p\equiv1$ or $5\pmod6$, and $1\le m<N$.
For any $c$ satisfying $c^p=v$, the difference-of-powers identity gives
\[
 (X-u)^m-c(X-v)^m
 \ \mid\ (X-u)^{pm}-v(X-v)^{pm}
 \qquad\text{in }K[X].
\]
If $p\equiv1\pmod6$, choose $c=v$.  Equation~\eqref{eq:a327315-closed}
then gives $S_m\mid S_N$ over $K$.

If $p\equiv5\pmod6$, choose $c=u$, since $u^5=v$.  Substituting $1-X$
in the closed form gives the exact reflected identity
\[
 S_m(1-X)=-u\big((X-u)^m-u(X-v)^m\big).
\]
Consequently $S_m(1-X)\mid S_N$ over $K$.

Both proposed divisors belong to $\mathbb Q[X]$.  Divisibility descends
to $\mathbb Q[X]$ by polynomial division and injectivity of the
coefficient embedding.  Each divisor has degree $m$, strictly between
zero and $N$, so $S_N$ is reducible.

Finally, a positive integer with no prime divisor other than 2 and 3
has the form $2^a3^b$, by induction on its prime factors.  This proves
the converse for every positive degree $N$.  Taking $N=n+1$ in
\eqref{eq:a327315-source} proves Theorem~\ref{thm:a327315}, including
the initial linear row.  Multiplication by the nonzero constant $n!$
does not change irreducibility.

\subsection{Prior literature and formal scope}

The source polynomial construction has an exact antecedent in
L.~Remete's \emph{A generalization of simplest number fields and their
integral basis}, arXiv:2111.08341v1 (2021)~\cite{Remete}.  His Section 2 defines
\[
 f_m^{(N)}(Y)=\sum_{j=0}^N\binom NjY^j g(N-j),
\]
where $g$ is six-periodic with values
$(1,-m,-m-1,-1,m,m+1)$.  Comparison with the definitions gives
$S_N(X)=(-1)^{N+1}f_1^{(N)}(X-1)$.
Indeed, at $m=1$, $g(k)=u^{k+1}+v^{k+1}$, so the binomial theorem gives
$f_1^{(N)}(Y)=u(Y+u)^N+v(Y+v)^N$, which agrees with the closed form
after the stated shift and sign.
The family itself, its derivative structure, and fractional-linear
root symmetries are therefore not new constructions here.
Remete's Corollary 10 requires a prime $p\ne3$ with
$v_p(m^2+m+1)=1$.  At $m=1$ that condition fails, and the inspected
irreducibility subsection does not state the complete degree criterion.

Related classical methods occur in C.~Cass, J.~Cullinan,
A.~Rasmussen and D.~Trifunovski, \emph{Specializations of Generalized
Rikuna Polynomials}, International Journal of Number Theory 10 (2014),
585--600~\cite{Rikuna}, and J.~S.~Calcut, \emph{The Tangent Analogues of the Chebyshev
Polynomials} (2008)~\cite{Calcut}.  Their relevant definitions and irreducibility or
angle-factorization sections were checked.  The present proof uses
standard cyclotomic irreducibility and fractional-linear substitution;
it does not claim those techniques as new.

The checked A327315 export still labels the exact classification a
conjecture.  The primary-reference review, exact and semantic web
queries, trureturing catalog, and targeted global and ProveIt GitHub
code searches located no earlier proof of that exact classification.
This is a bounded review, not a guarantee of historical priority.
The companion source-check report preserves the immutable OEIS exports,
their hashes, the index convention, and the scope of the prior-art search.

The formal development has nine local Lean modules.  It proves the
recurrence identities, exact degree, quadratic-field relations,
cyclotomic identity, both divisibility branches, the all-degree
criterion, and the exact canonical numerator bridge.  It uses the pinned Lean and mathlib versions
listed in Section~\ref{sec:verification}.
The rational-function derivative and fractional-linear helper are
reused unchanged from earlier packages.  The theorem dependencies are
the standard \texttt{propext}, \texttt{Classical.choice}, and
\texttt{Quot.sound}; no admitted proof or custom axiom is used.
